\begin{table*}[t]
\caption{\revisionA{Representative circuit templates and example compositions. }}
\label{tbl:comb}
  % \vspace{-0.2cm}
\centering
\scriptsize
\setlength{\tabcolsep}{3pt}
\renewcommand{\arraystretch}{1.35}

\begin{tabularx}{\textwidth}{@{} P{1.6cm} P{2.4cm} P{2.2cm} X P{3.2cm} @{}}
\toprule
\textbf{Current template} & \textbf{Full circuit template name} & \textbf{Next possible templates} & \textbf{Description} & \textbf{Example compositions} \\
\midrule

% ---- START ----
\famSTART
& (start symbol)
& \famStatePrep\candsep \famOracleA\candsep \famQFT\candsep \famVariational
& Common entry points (prepare input/eigenstate/ansatz), enabling realistic multi-block pipelines.
& \vis{\famSTART\qarrow\famStatePrep\\[-1pt]
       \famSTART\qarrow\famOracleA\\[-1pt]
       \famSTART\qarrow\famQFT} \\
\midrule % <-- horizontal line after START row

% ---- StatePrep ----
\famStatePrep
& State Preparation
& \famHamSim\candsep \famModularExp\candsep \famGroverIter\candsep \famVariational
& Prepares \(|\psi\rangle\)/\(|b\rangle\) or initial superpositions before an algorithmic primitive.
& \vis{\famStatePrep\qarrow\famHamSim\\[-1pt]
       % \famStatePrep\qarrow\famModularExp\qarrow\famQPE} \\
       \famStatePrep\qarrow\famVariational} \\
% ---- HamSim ----
\famHamSim
& Hamiltonian Simulation
& \famQPE\candsep \famVariational\candsep \famRandom
& Controlled time-evolution unitaries are typically read out via QPE (incl.\ IQFT decoding) in eigenvalue/spectral workflows.
& \vis{\famHamSim\qarrow\famQPE} \\       
% ---- QFT ----
\famQFT
& Quantum Fourier Transform
& \famQPE\candsep \famMeasure\candsep \famRandom
& Fourier-basis transform and dense benchmark block. 
& \vis{\famQFT\qarrow\famQPE\\[-1pt]
       \famQFT\qarrow\famMeasure} \\ 
% Variational 
% ---- Oracle/A ----
\famOracleA
& Oracle/Algorithm-\(A\) Wrapper 
& \famGroverIter\candsep \famVariational\candsep \famMeasure
& Encodes a problem-dependent subroutine \(A\) (e.g., oracle/feature map/encoding) that is naturally followed by Grover-style amplification.
% ; supports \(A \rightarrow\) GroverIter \(\rightarrow\) AmpEst.
& \vis{\Ablock\qarrow\famGroverIter\qarrow\famAmpEst} \\
% ---- ModularExp ----
\famModularExp
& Modular Exponentiation
& \famQPE
& Order-finding and Shor-style structure: modular arithmetic unitaries followed by phase estimation (incl.\ IQFT) to extract phases/periods.
% ; supports ModularExp $\rightarrow$ QPE.
& \vis{\famModularExp\qarrow\famQPE} \\
% ---- GroverIter ----
\famGroverIter
& Grover Iteration 
& \famGroverIter\ (repeat)\candsep \famAmpEst\candsep \famMeasure
& Amplification layers typically repeat and then transition to an estimation wrapper or terminate with measurement.
% ; supports \(A \rightarrow\) GroverIter \(\rightarrow\) AmpEst.
& \vis{\famGroverIter\qarrow\famGroverIter\\[-1pt]
       \famGroverIter\qarrow\famAmpEst\\[-1pt]
       \famGroverIter\qarrow\famMeasure} \\
% ---- AmpAmpl ----
\famAmpAmpl
& Amplitude Amplification
& \famMeasure
& Boosts the probability of a desired or postselected outcome. 
% e.g., using repeated Grover-style reflections. 
% In HHL-style compositions, this represents amplification of the successful postselection branch after Hamiltonian-simulation-driven phase estimation.
& \vis{\famAmpAmpl\qarrow\famMeasure} \\
% ---- QPE ----
\famQPE
& Quantum Phase Estimation
% & \famMeasure\candsep \famVariational\candsep \famRandom
&\famAmpAmpl\candsep \famMeasure\candsep \famVariational\candsep \famRandom
& Phase readout/decoding stage (typically includes an IQFT); often followed by measurement or optional continuations.
% & \vis{\famQPE\qarrow\famMeasure\\[-1pt]
%        \famQPE\qarrow\famVariational\\[-1pt]
%        \famQPE\qarrow\famRandom} \\
& \vis{\famQPE\qarrow\famAmpAmpl\\[-1pt]
       \famQPE\qarrow\famMeasure\\[-1pt]
       \famQPE\qarrow\famVariational\\[-1pt]
       \famQPE\qarrow\famRandom} \\


% ---- Variational ----
\famVariational
& Variational Circuit Layer/ Ansatz
& \famVariational\ (repeat)\candsep \famMeasure\candsep \famRandom
& Layered ansatz patterns; measurement used for objective estimation; can be interleaved for hybrid workloads.
& \vis{\famVariational\qarrow\famVariational\\[-1pt]
       \famVariational\qarrow\famMeasure\\[-1pt]
       \famVariational\qarrow\famRandom} \\

% ---- Random ----
\famRandom
& Random Circuit Block
& \famRandom\candsep \famMeasure
& Diversity/stress-test tail; may extend circuit length or terminate by measurement.
& \vis{\famRandom\qarrow\famRandom\\[-1pt]
       \famRandom\qarrow\famMeasure} \\

% ---- AmpEst ----
% \famAmpEst
% & Amplitude Estimation
% & ---, \famRandom %\candsep \famRandom
% & Estimation stage that is typically the final step, or proceeds to auxiliary stress-test blocks.
% & \vis{\famAmpEst\qarrow\famRandom}\\
% % \\[-1pt]
% %        \famAmpEst\qarrow\famRandom} 
% ---- AmpEst ----
\famAmpEst
& Amplitude Estimation
& \famMeasure
& Estimates a success amplitude or probability. It normally terminates the quantum subroutine by measurement and classical post-processing.
% ; it is distinct from amplitude amplification, which boosts a success probability.
& \vis{\famAmpEst\qarrow\famMeasure} \\
\midrule % <-- horizontal line after Modular Exponentiation row    
% ---- Measure ----

% & Measurement  
% & ---
% & Final step that yields classical outcomes and ends the whole circuit generation process.
% & \vis{\famMeasure} \\

% ---- Measure ----
\famMeasure
& Measurement  
& \famEND
& Final step that yields classical outcomes and ends the whole circuit generation process.
% Final emitted circuit/readout template. It yields classical outcomes, probabilities, or observable estimates, after which the generator transitions to the terminal END symbol.
& \vis{\famMeasure\qarrow\famEND} \\
% ---- END ----
\famEND
& (end symbol)
& ---
& 
% Absorbing generator sentinel that emits no gates or subcircuits. 
It marks the end of the template sequence.
& \vis{\famEND} \\
\bottomrule
\end{tabularx}


\end{table*}
